% =============================================================================
% APPENDIX AU: CORE-AWARE: The Awareness Function A(·)
% =============================================================================
% Module: BCM2_11_AWA (Awareness)
% Category: CORE
% Version: 1.0 (January 2026)
% Status: Final
% Language: English
%
% PURPOSE: This appendix answers the question "How aware is the individual of
% their options and consequences?" It formalizes the Awareness Function A(·)
% that modulates effective utility.
%
% 10C POSITION: This is CORE 6 of 10 - the "AWARE" question.
%
% DETAILED AXIOM SYSTEM: See Appendix AWA for the complete 80-axiom system.
%
% =============================================================================

\begin{tcolorbox}[colback=gray!5!white, colframe=gray!75!black,
    title=Appendix Header: AU CORE-AWARE]

\begin{tabular}{@{}ll@{}}
\textbf{Code:} & AU \\
\textbf{Category:} & CORE (Fundamental Theory) \\
\textbf{Full Name:} & CORE-AWARE: The Awareness Function \\
\textbf{Module:} & BCM2\_11\_AWA (Awareness) \\
\textbf{Version:} & 1.0 (January 2026) \\
\textbf{Status:} & Final \\
\end{tabular}

\vspace{0.5em}
\textbf{Purpose:} Answers the sixth fundamental question:
\textit{``How aware is the individual of their options, consequences, and current state?''}
Formalizes the Awareness Function $A(\cdot) \in [0, 1]$ that modulates effective utility.

\textbf{Detailed Treatment:} See Appendix AWA for the complete 80-axiom system.

\end{tcolorbox}

\begin{tcolorbox}[colback=purple!5!white, colframe=purple!75!black,
    title=Chapter Linkage]

\textbf{Primary Chapter:} Chapter 11: Awareness and Attention

\textbf{The Fundamental Question:}
\begin{quote}
\textit{``How aware is the decision-maker of their options, the consequences of their
choices, and their own preferences? How does limited awareness affect behavior?''}
\end{quote}

\textbf{Relationship to Other COREs:}
\begin{itemize}[nosep]
\item Modifies effective utility: $U_{\text{eff}} = A(\cdot) \times U$
\item Context-dependent: $A = A(\Psi)$ via CORE-WHEN (V)
\item Level-specific: $A_L$ varies by welfare level (CORE-WHO)
\end{itemize}

\end{tcolorbox}

\chapter{CORE-AWARE: The Awareness Function}
\label{app:core-aware}

\begin{abstract}
This appendix answers the sixth fundamental question:
\textit{``How aware is the individual?''}
We establish that behavior depends not just on utility ($U$) but on \textit{perceived}
utility, filtered through awareness: $U_{\text{eff}} = A(\cdot) \times U$.
The Awareness Function $A(\cdot) \in [0, 1]$ captures information availability,
attention allocation, salience, and metacognitive accuracy.
\end{abstract}

\begin{tcolorbox}[colback=blue!5!white, colframe=blue!75!black,
    title=Quick Reference: CORE-AWARE]

\textbf{Core Question:} How aware is the decision-maker?

\textbf{Key Formula:}
\begin{equation}
U_{\text{eff}} = A(\cdot) \cdot U \quad \text{where } A(\cdot) \in [0, 1]
\label{eq:aware-effective}
\end{equation}

\textbf{Awareness Components:}
\begin{itemize}[nosep]
\item $A_{\text{info}}$: Information availability (knows options exist)
\item $A_{\text{att}}$: Attention allocation (focuses on relevant factors)
\item $A_{\text{sal}}$: Salience (consequences are vivid)
\item $A_{\text{meta}}$: Metacognition (knows what they know/don't know)
\item $A_{\text{moral}}$: Moral awareness (perceives ethical implications, shaped by narratives)
\end{itemize}

\textbf{Cross-References:}
AWA (complete axiom system), V (CORE-WHEN), CAL (METHOD-LLMMC), CP (DOMAIN-CPS: Cognitive Load in Complex Problem Solving), RB (LIT-BENABOU: Moral Narratives)

\end{tcolorbox}

\section{The Fundamental Question}
\label{sec:aware-fundamental}

\begin{tcolorbox}[colback=red!5!white, colframe=red!75!black,
    title=The Fundamental Question]
\textbf{How aware is the decision-maker of their options, consequences, and
own preferences? How does limited or asymmetric awareness affect choices?}
\end{tcolorbox}

Standard economics assumes perfect information. Behavioral economics relaxes this
but often treats awareness as binary (know/don't know). We formalize awareness
as a continuous function with multiple components.

\section{Core Theory: The Awareness Function}
\label{sec:aware-theory}

\subsection{Definition}

\begin{definition}[Awareness Function]
The \textbf{Awareness Function} $A: \Omega \to [0, 1]$ maps states of the world
to awareness levels, where:
\begin{itemize}
\item $A = 0$: Complete unawareness (option/consequence not perceived)
\item $A = 1$: Full awareness (perfect perception)
\item $0 < A < 1$: Partial awareness (attenuated perception)
\end{itemize}
\end{definition}

\subsection{Components of Awareness}

\begin{equation}
A(\cdot) = f(A_{\text{info}}, A_{\text{att}}, A_{\text{sal}}, A_{\text{meta}}, A_{\text{moral}})
\end{equation}

\begin{table}[htbp]
\centering
\caption{Awareness Components}
\label{tab:aware-components}
\renewcommand{\arraystretch}{1.3}
\begin{tabular}{@{}clp{5cm}@{}}
\toprule
\textbf{Component} & \textbf{Name} & \textbf{Definition} \\
\midrule
$A_{\text{info}}$ & Information & Does the option/consequence exist in awareness? \\
$A_{\text{att}}$ & Attention & Is attention allocated to this factor? \\
$A_{\text{sal}}$ & Salience & How vivid/prominent is the factor? \\
$A_{\text{meta}}$ & Metacognition & Awareness of own knowledge state \\
$A_{\text{moral}}$ & Moral Awareness & Perception of ethical implications (narrative-dependent) \\
\bottomrule
\end{tabular}
\end{table}

\section{Axiom System}
\label{sec:aware-axioms}

\begin{tcolorbox}[colback=gray!5!white, colframe=gray!75!black,
    title=Axiom AWX-A1: Core Transformation]
\textbf{Statement:}
\begin{equation}
U_{\text{eff}} = A(\cdot) \cdot U_{\text{objective}}
\end{equation}

\textbf{Interpretation:} Effective utility equals objective utility filtered through awareness.
Unawareness of benefits makes them behaviorally irrelevant.
\end{tcolorbox}

\begin{tcolorbox}[colback=gray!5!white, colframe=gray!75!black,
    title=Axiom AWX-A2: Boundedness]
\textbf{Statement:}
\begin{equation}
A(\cdot) \in [0, 1] \quad \text{(normalized)}
\end{equation}

\textbf{Interpretation:} Awareness is bounded between complete unawareness (0) and
full awareness (1).
\end{tcolorbox}

\begin{tcolorbox}[colback=gray!5!white, colframe=gray!75!black,
    title=Axiom AWX-A3: Context Dependence]
\textbf{Statement:}
\begin{equation}
A = A(\Psi) \quad \text{where } \Psi \text{ is the context vector}
\end{equation}

\textbf{Interpretation:} Awareness is context-dependent. The same person has
different awareness levels in different situations.

\textbf{Cross-Reference:} CORE-WHEN (V) specifies $\Psi$.
\end{tcolorbox}

\begin{tcolorbox}[colback=gray!5!white, colframe=gray!75!black,
    title=Axiom AWX-A4: Asymmetry (Gains vs. Losses)]
\textbf{Statement:}
\begin{equation}
A_{\text{loss}} > A_{\text{gain}} \quad \text{(losses are more salient)}
\end{equation}

\textbf{Interpretation:} Negative outcomes capture attention more effectively
than positive outcomes of equivalent magnitude.
\end{tcolorbox}

\textbf{Note:} The complete axiom system (AWX-A1 through AWX-A80) is documented
in Appendix AWA. This appendix provides the core theoretical framework.

\section{Awareness Dynamics}
\label{sec:aware-dynamics}

\subsection{Decay}

\begin{equation}
A(t) = A_0 \cdot e^{-\kappa t} + A_{\text{stable}}
\end{equation}

where $\kappa$ is the decay rate and $A_{\text{stable}}$ is the baseline level.

\subsection{Trigger Reinforcement}

\begin{equation}
A_{\text{post-trigger}} = A_{\text{pre}} + \Delta A_{\text{trigger}} \cdot (1 - A_{\text{pre}})
\end{equation}

Repeated triggers increase baseline awareness (habituation with ceiling effect).

\section{Moral Awareness and Narratives}
\label{sec:aware-moral-narratives}

\begin{tcolorbox}[colback=orange!5!white, colframe=orange!75!black,
    title=Moral Awareness: How Narratives Shape Perception of Consequences]

Moral behavior depends critically on awareness of moral implications. Bénabou, Falk, and Tirole \citep{benabou_2018_narratives} show that \textbf{narratives} systematically manipulate this awareness:

\begin{equation}
A_{\text{moral}}(e) = A_{\text{base}} + \Delta A_{\text{narrative}}
\label{eq:moral-awareness}
\end{equation}

where $e$ is the perceived externality and $\Delta A_{\text{narrative}}$ can be positive (responsibilizing) or negative (exculpatory).

\end{tcolorbox}

\subsection{Exculpatory vs. Responsibilizing Narratives}

\begin{definition}[Moral Narratives]
\textbf{Exculpatory narratives} reduce perceived moral costs by:
\begin{itemize}[nosep]
\item Downplaying externalities: $e^* < e$ (``It doesn't really matter'')
\item Magnifying action costs: $c^* > c$ (``It's too hard to act'')
\item Denying pivotality: ``My contribution won't make a difference''
\end{itemize}

\textbf{Responsibilizing narratives} increase perceived moral costs by:
\begin{itemize}[nosep]
\item Highlighting externalities: $e^* > e$ (``Every action counts'')
\item Emphasizing agency: ``You can make a difference''
\item Social proof: ``Others are doing their part''
\end{itemize}
\end{definition}

\subsection{Network Transmission of Awareness}
\label{subsec:aware-network}

Moral awareness spreads through social networks with parameter $\rho$ (segregation):

\begin{equation}
A_{\text{moral},i}(t+1) = (1-\rho) \cdot \bar{A}_{\text{moral}} + \rho \cdot A_{\text{moral},\text{neighbors}}
\end{equation}

\begin{tcolorbox}[colback=gray!5!white, colframe=gray!75!black,
    title=Axiom AWX-A5: Narrative Transmission]
\textbf{Statement:}
\begin{itemize}[nosep]
\item Negative (exculpatory) narratives are \textbf{strategic substitutes}: If others adopt them, incentive to adopt decreases
\item Positive (responsibilizing) narratives are \textbf{strategic complements}: If others adopt them, incentive to adopt increases
\end{itemize}

\textbf{Implication:} Prosocial norms can be self-sustaining (positive feedback), while moral erosion faces natural limits.

\textbf{Source:} \citet{benabou_2018_narratives}, Propositions 1-3.
\end{tcolorbox}

\subsection{Image Concerns and Moral Awareness}

The utility function includes image concerns that interact with awareness:

\begin{equation}
U = (ve - c/\beta)a + \mu \hat{v}(a)
\end{equation}

where:
\begin{itemize}[nosep]
\item $\mu$ = strength of image concerns (self and social)
\item $\hat{v}(a)$ = inferred moral type from action $a$
\item $\beta$ = self-control parameter (hyperbolicity)
\end{itemize}

Higher $\mu$ increases incentive to \textit{seek} responsibilizing narratives (to justify costly moral action) or \textit{adopt} exculpatory narratives (to justify inaction while maintaining self-image).

\textbf{Cross-Reference:} See LIT-BENABOU (RB), Axiom RB-7 for the complete formalization.

\section{Integration with Other COREs}

\begin{tcolorbox}[colback=yellow!5!white, colframe=yellow!75!black,
    title=Integration: CORE-AWARE $\leftrightarrow$ Other COREs]
\textbf{With CORE-WHO (AAA):}
Awareness exists at each welfare level; group awareness aggregates individual awareness.

\textbf{With CORE-WHAT (C):}
Awareness varies by utility dimension; some dimensions are more salient than others.

\textbf{With CORE-WHEN (V):}
Context determines awareness levels: cognitive load reduces $A$, salience increases $A$.

\textbf{With CORE-READY (AV):}
Awareness is a prerequisite for willingness: $W > 0 \Rightarrow A > 0$.

\textbf{Application Domain: Complex Problem Solving (CP):}
Cognitive Load Theory (Sweller) maps directly to AWARE: $A(\Psi_{\text{load}}) = A_{\max} \times \exp(-k \times \Psi_{\text{load}})$. See Appendix CP for the full CPS-10C mapping.
\end{tcolorbox}

\section{Summary}

\begin{tcolorbox}[colback=gray!5!white, colframe=gray!75!black,
    title=Appendix AU Summary: CORE-AWARE]

\textbf{Core Contribution:}
Formalizes the Awareness Function $A(\cdot) \in [0, 1]$ that filters objective
utility into effective utility.

\textbf{Key Formula:}
\begin{equation}
U_{\text{eff}} = A(\cdot) \cdot U
\end{equation}

\textbf{Key Insights:}
\begin{itemize}[nosep]
\item Unawareness of options makes them behaviorally irrelevant
\item Awareness has four components: information, attention, salience, metacognition
\item Awareness is context-dependent and subject to decay/reinforcement
\item Losses are more salient than gains (asymmetric awareness)
\item \textbf{Moral narratives} can increase or decrease awareness of moral implications
\item Narratives spread through networks: positive = complements, negative = substitutes
\end{itemize}

\textbf{Detailed Treatment:} See Appendix AWA for complete 80-axiom system.

\end{tcolorbox}

\vspace{1em}
\hrule
\vspace{0.5em}

\noindent\textit{Cross-references:}
Appendix AWA (detailed axioms), V (CORE-WHEN), AV (CORE-READY), CAL (METHOD-LLMMC), RB (LIT-BENABOU: Moral Narratives), FAL (LIT-FALK), JT (LIT-TIROLE).

\noindent\textit{Master Bibliography:} \texttt{bcm\_master.bib}

\nocite{bcm_master}

% =============================================================================
% END OF APPENDIX AU
% =============================================================================
